\section{Architectural Trends and Related Work}

Deep neural networks have catalyzed significant progress in medical computer vision, yet their deployment in real-time neuro-oncology remains severely hindered by the compute, memory, and latency bottlenecks of prevailing architectures. In this section, we review relevant literature across dominant architectural paradigms, examining their mathematical foundations, adoption in brain MRI diagnosis, and intrinsic operational limitations on resource-constrained clinical hardware.

\subsection{Increasing Depth with Residuals in Brain MRI}
As convolutional neural networks grew deeper, optimization became increasingly hindered by vanishing and exploding gradients, as well as degradation where accuracy saturates and abruptly drops. He \emph{et~al.} \cite{14} resolved this degradation phenomenon with the Residual Network (ResNet), introducing identity shortcut connections that allow forward activations and backward error signals to bypass parameterized convolutional layers without attenuation. In clinical neuro-oncology, deep residual architectures---predominantly ResNet-18, ResNet-50, and ResNet-101---have established themselves as the standard backbone across tumor detection, multi-class classification, and histologic glioma grading \cite{10, 38, 54, 55}. Extensive surveys by Litjens \emph{et~al.} \cite{10}, Lundervold and Lundervold \cite{54}, and Nayak \emph{et~al.} \cite{55} demonstrate that transfer learning from ImageNet-pretrained ResNet-50 models yields near-saturated diagnostic accuracies exceeding 98\% on standard multi-class MRI benchmarks, such as the Nickparvar dataset \cite{34} and Figshare.

Despite their high benchmark accuracy, deep residual networks introduce acute operational liabilities in clinical practice. In original architectural specifications and comprehensive benchmarks by He \emph{et~al.} \cite{14} and Bianco \emph{et~al.} \cite{60}, standard ResNet-50 contains approximately 25.56 million parameters, requiring roughly 98 to 102 megabytes of storage for uncompressed single-precision weights alone, while VGG-16 \cite{39} comprises upwards of 138 million parameters (and over 27 million even when dense classification heads are trimmed). In terms of arithmetic complexity, a standard single-image forward pass through ResNet-50 requires 3.8 to 4.1 GFLOPs (roughly 3.86 to 4.09 GMACs) \cite{14, 60}, whereas VGG-16 demands over 15.3 GMACs (exceeding 30 GFLOPs) \cite{39, 60}. As documented in comprehensive hardware surveys by Mittal \emph{et~al.} \cite{33} and Amin \emph{et~al.} \cite{56}, hospital workstations and portable intra-operative scanners rely primarily on general-purpose central processing units (CPUs) or low-power embedded processors lacking high-end enterprise GPU acceleration. On such resource-constrained hardware, deep residual backbones encounter severe latency bottlenecks: empirical benchmarking across embedded and edge platforms by Bianco \emph{et~al.} \cite{60} reveals that ResNet-50 is throttled to under 19 frames per second, with single-image inference latency frequently exceeding 50 to 80 milliseconds on unaccelerated consumer CPUs \cite{33, 60, 61}. For clinical 3D MRI examinations containing hundreds of volumetric slices across multiple modalities (T1, T1Gd, T2, FLAIR) \cite{2, 28}, this computational overhead translates into cumulative multi-second delays per patient, impeding emergency neuro-triage and interactive intra-operative navigation. Furthermore, as demonstrated by Ma \emph{et~al.} \cite{32}, the multi-branch bottleneck topology of ResNets incurs substantial Memory Access Cost (MAC) and off-chip DRAM memory traffic, generating severe cache thrashing that causes real-world hardware execution time to deviate unfavorably from theoretical FLOP estimates.

\subsection{The Depthwise Separable Convolution Paradigm}
To provide computationally lightweight models capable of real-time execution on constrained hardware, Howard \emph{et~al.} \cite{15} proposed MobileNet, replacing standard spatial convolutions with depthwise separable convolutions. This decomposition decouples spatial filtering from channel mixing: a depthwise convolution applies a single spatial kernel per input channel, followed by a $1\times1$ pointwise convolution that linearly combines features across channels. For $3\times3$ kernels, this factorization achieves an 8- to 9-fold reduction in arithmetic operations with minimal loss of representational capacity. Sandler \emph{et~al.} \cite{37} introduced MobileNetV2, pioneering the inverted residual block with a linear bottleneck. Unlike conventional residual blocks that compress feature dimensions before spatial convolution, MobileNetV2 expands low-dimensional input features to a higher-dimensional intermediate space using pointwise convolutions, performs spatial depthwise filtering, and subsequently projects features back to a lower-dimensional manifold via a linear bottleneck to prevent non-linear activation collapse. MobileNetV3 \cite{19} enhanced this framework by pairing neural architecture search (NAS) \cite{50} with hardware-aware optimization, integrating hard-swish non-linearities and selective Squeeze-and-Excitation (SE) channel recalibration blocks \cite{16} to dynamically reweight feature maps according to global context.

\subsection{Compound Scaling}
The EfficientNet family \cite{40} was formulated to systematically balance accuracy and computational demand by rethinking network scaling. Prior works typically scaled a single dimension---either network depth, channel width, or input image resolution---in isolation. Tan and Le demonstrated that these three dimensions are fundamentally interdependent and introduced a principled compound scaling method that uniformly scales depth, width, and resolution using a fixed compound coefficient, effectively constraining target multiply-accumulate operations (MACs). EfficientNet-B0 serves as the baseline, operating at 5.3 million parameters and 0.39 GMACs \cite{40}. To resolve training latency bottlenecks and reduce memory overhead, EfficientNetV2 \cite{41} incorporated Fused-MBConv blocks, replacing early depthwise separable blocks with standard $3\times3$ convolutions to better exploit hardware acceleration and memory locality on modern CPUs and GPUs, while employing progressive learning strategies.

\subsection{Shuffling and Squeezing}
Aiming directly at hardware runtime rather than indirect FLOP metrics, Ma \emph{et~al.} \cite{32} developed ShuffleNetV2 based on four hardware-aware design principles: maintaining equal input and output channel widths to minimize memory access cost, minimizing excessive group convolutions, reducing network branching, and eliminating element-wise additions. ShuffleNetV2 divides input channels into two paths via a channel split layer, conducts depthwise convolutions along one path, concatenates the resulting branches, and applies a channel shuffle operation to facilitate cross-group feature exchange without matrix multiplication. Concurrently, SqueezeNet \cite{17} demonstrated that parameter counts could be aggressively reduced through ``Fire'' modules, which first squeeze feature channels via $1\times1$ convolutions and then expand spatial representations through parallel $1\times1$ and $3\times3$ convolutions, matching AlexNet-level recognition while maintaining an uncompressed footprint below 1.25 million parameters. GhostNet \cite{13} further reduced computational redundancy by observing that feature maps often contain mutual similarities; it generates a compact set of intrinsic feature maps via standard convolution and synthesizes remaining ``ghost'' features through computationally inexpensive linear depthwise operations.

\subsection{Attention and Modernized ConvNets}
Moving beyond pure convolutional inductive biases, the Vision Transformer (ViT) \cite{8} tokenizes images into discrete coordinate patches and applies global self-attention mechanisms \cite{43}. While ViTs exhibit exceptional representational capacity when pre-trained on web-scale datasets, their lack of translational equivariance and local inductive priors makes them notoriously prone to severe overfitting when trained from scratch on small-to-medium-scale medical cohorts \cite{24}. Furthermore, the quadratic computational complexity of pairwise global self-attention with respect to token length imposes prohibitive latency and memory footprints during high-resolution volumetric MRI processing \cite{26}. Motivated by transformer advancements, Liu \emph{et~al.} \cite{30} introduced ConvNeXt, modernizing classic ResNets by incorporating large $7\times7$ kernels, inverted bottleneck designs, depthwise convolutions, and LayerNorm to achieve transformer-competitive performance while retaining the hardware efficiency and simplicity of ConvNets. ConvNeXt V2 \cite{46} combined this architecture with masked autoencoders and global response normalization to stabilize feature representation learning.

\subsection{Multi-Scale Receptive Fields and Atrous Convolutions}
Discriminating diffuse intracranial tumor boundaries, perifocal vasogenic edema, and necrotic cores requires capturing spatial context across heterogeneous scales \cite{2, 28, 55}. Standard CNNs capture multi-scale context by progressively downsampling feature maps, sacrificing spatial boundary precision. In contrast, atrous (dilated) convolutions \cite{5} expand the effective receptive field without spatial resolution reduction or parameter inflation by introducing an integer dilation rate that spaces out kernel filter weights. Chen \emph{et~al.} \cite{5} popularized this concept through Atrous Spatial Pyramid Pooling (ASPP) in DeepLab, deploying parallel atrous convolutional branches with distinct dilation rates to capture multi-scale representations simultaneously. Kyrkou and Theocharides \cite{22} utilized multi-scale atrous fusion in EmergencyNet to achieve real-time aerial classification. Building upon these principles, Kyrkou \cite{21} introduced LiteSpeed-Net (LS-Net), demonstrating that multi-path atrous convolutions provide a potent inductive bias for compact vision models operating under a sub-1M parameter budget: by splitting channels across parallel depthwise branches with dilation rates of 1, 2, and 3, LS-Net achieves broad multi-scale receptive field coverage while maintaining high inference speed. However, multi-path atrous convolutions suffer from an unaddressed vulnerability in clinical neuro-imaging: because dilated kernels sample the spatial grid at rigid, discrete intervals without spatial interpolation \cite{5, 21}, their predictions become disproportionately fragile under input scale variations and cross-institution acquisition shifts, such as varying slice thickness and lesion zoom levels.

\subsection{Test-Time Augmentation and Uncertainty Calibration}
In safety-critical clinical applications, deep models must communicate predictive uncertainty alongside raw class scores \cite{10, 57}. In neuro-oncology MRI, diagnostic ambiguity arises from anatomical heterogeneity, scanner noise, and indistinct lesion margins---collectively defined as data-dependent aleatoric uncertainty \cite{4, 47}. Seminal works by Ayhan and Berens \cite{4} and Wang \emph{et~al.} \cite{47} demonstrated that Test-Time Augmentation (TTA) provides an effective Monte Carlo sampling framework to estimate aleatoric uncertainty without requiring Bayesian retraining. By applying a distribution of stochastic geometric and photometric transformations to a query slice at test time and averaging the resulting predictions, TTA smooths the decision boundary and yields calibrated uncertainty estimates \cite{4, 47, 58}. Nevertheless, conventional TTA introduces a severe computational bottleneck for edge deployment: running an unselective multi-view ensemble (e.g., 6 to 8 forward passes per image) multiplies inference latency by an equivalent linear factor \cite{58}. In clinical edge workflows, executing an unconditional multi-view TTA scales latency linearly with the number of augmentation passes, completely negating the execution speedups attained through compact architecture design \cite{47, 58}. Furthermore, as shown by Shanmugam \emph{et~al.} \cite{58}, unconditionally augmenting unambiguous high-confidence samples can distort canonical anatomical orientations, inadvertently corrupting diagnostic accuracy.

\subsection{Main Limitations and Research Gaps}
A critical synthesis of existing literature exposes three fundamental voids that motivate this investigation. First, prior medical imaging studies focus almost exclusively either on overparameterized backbones exceeding 20 million parameters (ResNet-50, VGG-16) or mobile-scale models (2M to 20M parameters). The strict sub-1M parameter regime ($<1,000,000$ parameters)---which represents the critical threshold for low-power microcontroller deployment, zero-latency caching, and localized hospital workstation integration---remains unbenchmarked across in-domain classification and out-of-domain clinical grading. Second, while parallel atrous convolutions achieve high efficiency on saturated in-domain benchmarks, their discrete sampling mechanics induce acute vulnerability to cross-dataset scale shifts, where diagnostic errors concentrate in low-confidence predictions near lesion boundaries. Third, existing test-time augmentation pipelines operate as unselective, brute-force ensembles that impose prohibitive latency penalties. No prior methodology leverages predictive uncertainty to selectively trigger TTA while matching the test-time transformations to the internal multi-scale receptive fields of compact atrous architectures.

To resolve these limitations, this paper establishes a systematic sub-1M benchmark on brain MRI diagnosis and introduces Multi-Scale Consistency Gated Test-Time Augmentation (MCG-TTA), recovering out-of-domain clinical accuracy with zero additional learned parameters while preserving edge-level inference speed.
